\documentclass[a4paper,12pt]{article}

% Pacchetti fondamentali
\usepackage[T1]{fontenc}    % Codifica font (inputenc utf8 è predefinito dal 2018)
\usepackage[italian]{babel} % Lingua italiana
\usepackage{import}
% Colori (deve essere prima di pacchetti basati su tikz per evitare option clash)
\usepackage[table]{xcolor}  % Per colorare tabelle e testo

% Pacchetti matematici
\usepackage{mathtools}      % Pacchetto matematico essenziale (carica amsmath)
\usepackage{cancel}         % Per barrare termini nelle formule
% \usepackage{amssymb}        % Simboli matematici aggiuntivi (spesso in conflitto con newtxmath)

% Font e Teoremi
\usepackage{newtxtext}                     % Font Times per il testo
\usepackage{mleftright}                    % Per \mleft e \mright che funzionano bene con newtxmath
\usepackage{amsthm}         % Per teoremi e dimostrazioni (caricato prima di newtxmath per evitare conflitti con \openbox)
%\usepackage{newtxtext, newtxmath} % Font Times per testo e matematica
\usepackage[subscriptcorrection,slantedGreek,nofontinfo,
            eucal,   % \mathcal Euler Script
            eufrak,  % \mathfrak Euler Fraktur
            mtpscr,  % \mathscr Times-compatible (con lettere alternative)
            mtphrb,  % \mathbb Holey Roman Bold
            ]{mtpro2}
% \usepackage[upint]{newtxmath}              % Base matematica Times
% \usepackage[nonpxmath,scale=0.95]{eulerpx} % Lettere matematiche in AMS Euler scalate perfettamente
% \usepackage{libertine, libertinust1math}      % Font Linertine per testo e matematica

\usepackage{interval}
% Se preferisci le parentesi tonde (a, b) invece delle quadre rovesciate ]a, b[ per gli aperti, decommenta la riga sotto:
\intervalconfig{soft open fences}

\usepackage[spaced=-1]{diffcoeff} % Per differenziali e derivate con notazione intelligente

% Comandi per differenziali 
\difdef { l } { dn } { style = d^ } % Per differenziali di ordine superiore (es. d^2 x)
\NewDocumentCommand \dn { m m } { \dl.dn.[#1] #2 }

\difdef { l } { p }
{ op-symbol = \partial } % Per “differenziali parziali”
\NewDocumentCommand \dlp {} { \dl.p. }

% Comandi per Jacobiani
\difdef { j } { / } { style = / } % Per Jacobiani in line 

% Comandi per la valutazione di derivate in un punto specifico
\difdef { f, fp, s, sp } { | }
{
outer-Ldelim
= \left . ,
outer-Rdelim = \right |,
sub-nudge
= 0 mu
}
\difdef { f, fp, s, sp } { ] }
{
outer-Ldelim = \left [ ,
outer-Rdelim = \right ],
elbowroom = 1 mu,
sub-nudge = 0 mu
}

% Riduzione dello spazio prima della funzione (es. per f(x))
\difdef { f, fp } {} 
{
  derivand-sep = 2 mu plus 1 mu minus 1 mu % default era 3 mu
}

% 2. Riduzione dello spazio per le derivate in linea (slash fraction)
\difdef { s, sp } {} 
{
  derivand-sep = 1.3 mu plus 0.5 mu minus 0.5 mu % default era 2 mu
} 

% ==============================================================================
% ALFABETI CALLIGRAFICI
% ==============================================================================
% \mathcal è Euler Script (opzione eucal)
% \mathscr è Math Script (opzione mtpscr) con supporto ad \altC, \altG, ecc.

% Curly Font (Times-compatible upright script)
% Utilizziamo il wrapper interno \MTPsetupCurly per far funzionare le lettere 
% alternative specifiche per Curly (\altG, \altM, \altN, \altQ, \altY).
\DeclareRobustCommand*{\mathcurly}[1]{{\MTPsetupCurly\MTPCurly{#1}}}
% Nota: non definiamo \mathbcurly perché il font Curly non ha una variante grassetto.

% Computer Modern Calligraphic (il classico di default di LaTeX)
\DeclareMathAlphabet{\mathcmcal}{OMS}{cmsy}{m}{n}
\SetMathAlphabet{\mathcmcal}{bold}{OMS}{cmsy}{b}{n}
\DeclareMathAlphabet{\mathbcmcal}{OMS}{cmsy}{b}{n}

% ==============================================================================
% ALFABETI FRAKTUR
% ==============================================================================
% MathTime Professional Fraktur
% Utilizziamo il wrapper interno per far funzionare \altx, \alty e \altY
\DeclareRobustCommand*{\mathmtfrak}[1]{{\MTPsetupFrak\MTPFrak{#1}}}


% ==============================================================================
% ALFABETI BLACKBOARD BOLD (DOPPIO TRATTO)
% ==============================================================================
% \mathbb è Holey Roman Bold (opzione mtphrb)

% Holey Roman Bold Italic (Strutture generiche: K, A...)
\DeclareMathAlphabet{\mathbbi}{U}{mt2hrb}{m}{it}

% Holey Roman Bold Dark (Per versori i, j, k)
\DeclareMathAlphabet{\mathbbd}{U}{mt2hrb}{b}{n}

% MathTime Blackboard Bold (Dritto, pieno)
\DeclareMathAlphabet{\mathmtbb}{U}{mt2bb}{m}{n}

% MathTime Blackboard Bold Italic (Corsivo, pieno)
\DeclareMathAlphabet{\mathmtbbi}{U}{mt2bb}{m}{it}

% MathTime Blackboard Bold Dark (Grassetto, pieno)
\DeclareMathAlphabet{\mathmtbbd}{U}{mt2bb}{b}{n}

% Altri pacchetti
% \usepackage{thmtools}      % Per teoremi e dimostrazioni (alternativa ad amsthm)
\usepackage{enumitem}       % Per personalizzare gli elenchi
\usepackage{multirow}       % Per celle che si estendono su più righe
\usepackage{relsize}        % Per cambiare la dimensione del testo in modo relativo
\usepackage[b]{esvect}      % Per vettori con freccia più elegante
\usepackage{ifthen}         % PER I COMANDI CON LOGICA CONDIZIONALE
\usepackage{tikz-cd}        % Per diagrammi commutativi
\usepackage{graphicx}       % Per includere immagini
\usepackage{fancyvrb}       % Per ambienti verbatim non standard
\usepackage{fvextra}        % Estensioni per verbatim (es. \SaveVerb, \UseVerb)
\usepackage[Glenn]{fncychap}% Per capitoli con stile personalizzato 
%Options: Sonny, Lenny, Glenn, Conny, Rejne, Bjarne, Bjornstrup

\usepackage{geometry}       % Per gestire i margini
\geometry{a4paper, margin=2.5cm}

% setup intestazioni e piè di pagina
\usepackage{fancyhdr}
\pagestyle{fancy}
\fancyhf{} % Pulisce tutti gli header e footer
\setlength{\headwidth}{\textwidth} % Larghezza header uguale a textwidth/linewidth
\fancyhead[LE,RO]{\thepage} % Numeri di pagina ai lati esterni
\fancyhead[RE]{\nouppercase{\leftmark}} % Nome capitolo a sinistra nelle pagine pari
\fancyhead[LO]{\nouppercase{\rightmark}} % Nome sezione a destra nelle pagine dispari
\renewcommand{\headrulewidth}{0.4pt} % Linea sotto l'header
\renewcommand{\footrulewidth}{0pt} % Nessuna linea nel footer

% hyperref per link e riferimenti
\usepackage{hyperref}
\hypersetup{
    colorlinks=true,
    linkcolor=blue,
    filecolor=magenta,
    urlcolor=cyan,
    pdftitle={Lemma di Barbalat},
    pdfauthor={Davide De Marco},
    pdfsubject={},
    pdfkeywords={Lemma di Barbalat, Stabilità, Teoria del Controllo, Lyapunov, Continuità Uniforme, derivate, convergenza},
}

% Comando intelligente per costanti matematiche
\ExplSyntaxOn
\NewDocumentCommand{\upconst}{m}
 {
  % Controlla se il numero di caratteri dell'input (#1) è esattamente 1
  \int_compare:nNnTF { \str_count:n { #1 } } = { 1 }
   { \mathrm { #1 } }    % [Vero] 1 carattere: applica \mathrm
   { \use:c { up #1 } }  % [Falso] >1 caratteri: equivalente a \csname up#1\endcsname
 }
\ExplSyntaxOff

% Costanti matematiche e versori
\newcommand{\eu}{\upconst{e}}           % Costante di Nepero
\newcommand{\iu}{\upconst{i}}           % Unità immaginaria
\newcommand{\piu}{\upconst{pi}}               % Pi greco
\newcommand{\vsi}{\mathbbd{i}}     % Versore asse x
\newcommand{\vsj}{\mathbbd{j}}     % Versore asse y
\newcommand{\vsk}{\mathbbd{k}}          % Versore asse z

% --- Definizioni Matematiche Personalizzate ---
\newcommand{\K}{\mathbbi{K}}             % Campo K

% Operatori
\DeclareMathOperator{\rk}{rk}           % Rango
\DeclareMathOperator{\Sol}{Sol}         % Spazio delle Soluzioni
\DeclareMathOperator{\GL}{GL}           % Gruppo Lineare Generale
\DeclareMathOperator{\Imm}{im}          % Insieme immagine
\DeclareMathOperator{\Span}{span}       % Span
\DeclareMathOperator{\Spanbis}{\mathscr{\altL}} % Span alternativo con L calligrafica
\DeclareMathOperator{\Cof}{cof}         % Cofattore
\DeclareMathOperator{\Adj}{adj}         % Matrice aggiunta (adjugate)

% Operatori differenziali
\newcommand{\grad}{\nabla\!} % Gradiente con spazio negativo per avvicinarlo alla funzione
\newcommand{\diverg}{\nabla \cdot}
\newcommand{\rot}{\nabla \times}
\newcommand{\lapl}{\nabla^2} % Oppure \Delta

% Matrici speciali
\newcommand{\Id}{\mathbbd{1}}           % Matrice Identità (blackboard bold)

%classe C di funzioni
\newcommand{\classC}[1]{\mathscr{\altC}^{#1}} % Classe C^k di funzioni

% Spazio delle Matrici "Intelligente"
% Sintassi: \Mat{righe}{colonne}{campo}
\newcommand{\Mat}[3]{%
    \ifthenelse{\equal{#1}{#2}}         % Controlla se righe=colonne
        {\mathscr{M}_{#1}(#3)}          % Se righe=colonne
        {\mathscr{M}_{#1,#2}(#3)}       % Se righe!=colonne
}

% Trasposta, Vettori e Matrici
\newcommand{\vect}[1]{\mathbold{#1}}   % Vettore 
\newcommand{\rvect}[1]{\vv{#1}} % Raggio vettore con freccia
\newcommand{\tr}{^{\normalfont\textsf{\relsize{-1}T}}}               % Trasposta
\newcommand{\matr}[1]{\mathbold{#1}} % Matrice (grassetto)

% Barra di separazione otticamente alleggerita per matrici a blocchi
\newcommand{\middledivbar}{\;\color{black!46}\middle|\color{black}\;} % Per matrici a blocchi adattabile
\newcommand{\bigdivbar}{\;\color{black!46}\big|\color{black}\;} % Per matrici a blocchi grandi
\newcommand{\Bigdivbar}{\;\color{black!46}\Big|\color{black}\;} % Per matrici a blocchi ancora più grandi

% --- NOTAZIONI MATRICIALI ---

% Cofattore di una matrice (con i,j come indici)
\newcommand{\cof}[3]{\Cof_{#1,#2}\mleft( \matr{#3} \mright)}

% 1. Elemento di posto i,j (Opzionale: potresti anche scriverlo a mano essendo molto semplice)
% Uso: \entry{A}{ij} oppure \entry{A}{i,j}
% \newcommand{\entry}[2]{#1_{#2}}

% 2. Riga di una matrice
% Uso: \matrow{A}{i}
\newcommand{\matrow}[2]{\matr{#1}_{(#2)}}

% 3. Colonna di una matrice
% Uso: \matcol{A}{j}
\newcommand{\matcol}[2]{\matr{#1}^{(#2)}}
% 4. Sottomatrice generale (estrazione di righe I e colonne J)
% Uso: \submat{A}{I}{J}
% Uso: \submat{A}{i_1,i_2,...,i_k}{j_1,j_2,...,j_l}
\newcommand{\submat}[3]{\matr{#1}_{(#2)}^{(#3)}}

% 5. Sottomatrice con riga i e colonna j eliminate (taglio)
% Uso: \subdel{A}{i}{j}
\newcommand{\subdel}[3]{\matr{#1}_{(#2 \mid #3)}}

\newcommand{\matrowsysm}[1]{\mathscr{R}_{\matr{#1}}} % Simbolo per il sistema di righe di una matrice
\newcommand{\matcolsysm}[1]{\mathscr{\altC}_{\matr{#1}}} % Simbolo per il sistema di colonne di una matrice
\newcommand{\req}{%
    \overset{%
      \text{%
        \raisebox{-0.4ex}[1ex][0ex]{\relsize{-2}$\mathscr{R}$}%
        }%
        }{\equiv}} % Riga-equivalenza con R sopra

% Insieme vuoto 
% \renewcommand{\emptyset}{\varnothing}   % Simbolo di insieme vuoto più elegante

\newcommand{\thmcita}[1]{\textsc{#1}} % Per citare teoremi con il nome in maiuscoletto


% Definizione di un array per matrici a blocchi con linee grigie
\newenvironment{blockarray}[1]
  {%
   \arrayrulecolor{black!46}% Colora le linee di grigio
   \begin{array}{#1}%
  }
  {%
   \end{array}%
   \arrayrulecolor{black}% Ripristina il nero per non intaccare tabelle future
  }

% Comandi per le operazioni elementari sulle righe
\newcounter{descrOP} % 1. Definisci il contatore

\makeatletter
\renewcommand{\p@descrOP}[1]{\normalfont OP\textsubscript{#1}}
\makeatother

% 2. Crea un comando personalizzato per l'item
\newcommand{\itemOp}{%
  \refstepcounter{descrOP}% Incrementa il contatore e permette a \label di vederlo
  \item[(OP\textsubscript{\thedescrOP})]% Stampa il numero seguito dal testo dell'opzione
}

\newcounter{descrDet} % 1. Definisci il contatore

\makeatletter
\renewcommand{\p@descrDet}[1]{\normalfont D\textsubscript{#1}}
\makeatother

\newcommand{\itemDet}{%
  \refstepcounter{descrDet}% Incrementa il contatore e permette a \label di vederlo
  \item[(D\textsubscript{\thedescrDet})]% Stampa il numero seguito dal testo dell'opzione
}

% Simbolo di somma con grandezza adattabile al contesto (funziona male con TX)
\makeatletter
\newcommand{\bigplus}{%
  \DOTSB\mathop{\mathpalette\mattos@bigplus\relax}\slimits@
}
\newcommand\mattos@bigplus[2]{%
  \vcenter{\hbox{%
    \sbox\z@{$#1\sum$}%
    \resizebox{!}{0.9\dimexpr\ht\z@+\dp\z@}{\raisebox{\depth}{$\m@th#1+$}}%
  }}%
  \vphantom{\sum}%
}
\makeatother

% --- Ambienti Teorema ---
\newtheoremstyle{normaltheorem}% Nome dello stile
    {3pt}% Space above
    {3pt}% Space below
    {\itshape}% Body font
    {}% Indent amount
    {\bfseries\scshape}% Theorem head font
    {.}% Punctuation after theorem head
    {.5em}% Space after theorem head
    {\thmname{#1}\thmnumber{ #2}\thmnote{ \normalfont\bfseries(#3)}}% Theorem head spec (can be left empty, meaning ‘normal’)
\theoremstyle{normaltheorem} % Stile corsivo per i teoremi (default)
\newtheorem{theorem}{Teorema}[section]
\newtheorem{lemma}[theorem]{Lemma}
\newtheorem{corollary}[theorem]{Corollario}
\newtheorem{proposition}[theorem]{Proposizione}

\newtheoremstyle{nametheorem}% Nome dello stile
    {3pt}% Space above
    {3pt}% Space below
    {\itshape}% Body font
    {}% Indent amount
    {\bfseries\scshape}% Theorem head font
    {.}% Punctuation after theorem head
    {.5em}% Space after theorem head
    {\thmname{#1}\thmnote{ #3}\thmnumber{ #2}}% Theorem head spec (can be left empty, meaning ‘normal’)
\theoremstyle{nametheorem}
\newtheorem{nametheorem}[theorem]{Teorema}
\newtheorem{namelemma}[theorem]{Lemma}

\makeatletter
\renewenvironment{proof}[1][\proofname]{\par
  \pushQED{\qed}%
  \normalfont \topsep6\p@\@plus6\p@\relax
  \trivlist
  \item\relax
  {\scshape
    #1\@addpunct{.}}\hspace\labelsep\ignorespaces
  }{%
    \popQED\endtrivlist\@endpefalse
  }
\makeatother

%\theoremstyle{definition} % Stile tondo per le definizioni
%\newtheorem*{definizione}{\normalfont\textsc{Definizione}}%[section]
\newtheoremstyle{definition_style}% Nome dello stile
    {3pt}% Space above
    {3pt}% Space below
    {}% Body font
    {}% Indent amount
    {\bfseries\scshape}% Theorem head font
    {.}% Punctuation after theorem head
    {.5em}% Space after theorem head
    {}% Theorem head spec (can be left empty, meaning ‘normal’)
\theoremstyle{definition_style}
\newtheorem{definition}{Definizione}[section]

\newtheoremstyle{notation_style}% Nome dello stile
    {3pt}% Space above
    {3pt}% Space below
    {}% Body font
    {}% Indent amount
    {\bfseries\scshape}% Theorem head font
    {.}% Punctuation after theorem head
    {.5em}% Space after theorem head
    {}% Theorem head spec (can be left empty, meaning ‘normal’)
\theoremstyle{notation_style}
\newtheorem{notation}[definition]{Notazione}

\newtheoremstyle{example_style}% Nome dello stile
    {3pt}% Space above
    {3pt}% Space below
    {}% Body font
    {}% Indent amount
    {\bfseries\scshape}% Theorem head font
    {.}% Punctuation after theorem head
    {.5em}% Space after theorem head
    {}% Theorem head spec (can be left empty, meaning ‘normal’)
\theoremstyle{example_style}
\newtheorem{example}{Esempio}[section]


\theoremstyle{remark} % Stile per le osservazioni
\newtheorem*{remark}{Osservazione}%[section]
\newtheorem*{property}{Proprietà}%[section]

%\setlist[description]{font=\normalfont\itshape, style=nextline}


\title{%
  \textbf{Lemma di Barbalat}\\
  %\rule[-1cm]{0pt}{1cm} % linea vuota per spaziatura
}
\author{%
  Davide De Marco
  %\\{\relsize{-1}Bacoli, Italia}
}
\date{}

\begin{document}
\maketitle
\section{Enunciato e dimostrazione del Lemma di Barbalat}
\begin{namelemma}[di Barbalat]\label{lemma:barbalat}
  Sia $f \, \colon \interval[open right]{0}{+\infty} \to \mathbb{R}$ una funzione tale che:
  \begin{enumerate}[label=\upshape({\itshape\roman*}\kern 0.6pt), ref=({\itshape\roman*}\kern 0.6pt)]
    \item $f \in \classC{1}\bigl(\interval[open right]{0}{+\infty}\bigr)$;
    \item $\lim\limits_{t \to + \infty} f(t)$ esiste e vale $\mathscr{l} \in \mathbb{R}$;
    \item $\dot{f}$ è uniformemente continua su $\interval[open right]{0}{+\infty}$.
  \end{enumerate}
  Allora, $\lim\limits_{t \to + \infty} \dot{f}(t) = 0$.
\end{namelemma}

\begin{proof}
  Si supponga per assurdo che 
  \[
  \lim_{t \to +\infty} \dot{f}(t) \neq 0
  \]
  oppure che $\dot{f}$ non ammetta limite per $t \to +\infty$. In entrambi i casi, negando la definizione di limite, esiste una costante $\varepsilon > 0$ e una successione $\mleft\{ t_n \mright\}_{n \in \mathbb{N}}$ tale che:
  \begin{equation}\label{eq:barbalat_contr}
    \lim_{n \to +\infty}t_n = +\infty \quad \text{e} \quad \mleft|\dot{f}(t_n) \mright| > \varepsilon \quad \forall n \in \mathbb{N}.
  \end{equation}
  Poiché $\dot{f}$ è uniformemente continua, per ogni $\bar{\varepsilon} > 0$ esiste $\delta = \delta(\bar{\varepsilon}) > 0$ tale che
  \begin{equation*}
    \mleft| \dot{f}(t) - \dot{f}(t_n) \mright| < \bar{\varepsilon} \quad \forall t \in \interval[open]{t_n}{t_n + \delta}, \,\forall n \in \mathbb{N}.
  \end{equation*}
  Scegliendo $\bar{\varepsilon} = \varepsilon/2$ si ottiene la relazione
  \begin{equation}\label{eq:barbalat_uc2}
    \mleft| \dot{f}(t) - \dot{f}(t_n) \mright| < \frac{\varepsilon}{2} \quad \forall t \in \interval[open]{t_n}{t_n + \delta}, \,\forall n \in \mathbb{N}.
  \end{equation}
  Mediante la disuguaglianza triangolare si ricava la stima
  \begin{equation*}
    \mleft| \dot{f}(t_n) \mright| = \mleft|\dot{f}(t) - \dot{f}(t_n) - \dot{f}(t) \mright| \leq \mleft| \dot{f}(t) - \dot{f}(t_n) \mright| + \mleft| \dot{f}(t) \mright|
  \end{equation*}
  da cui segue necessariamente
  \begin{equation}\label{eq:barbalat_triang2}
    \mleft| \dot{f}(t) \mright| \geq \mleft| \dot{f}(t_n) \mright| - \mleft| \dot{f}(t) - \dot{f}(t_n) \mright| \quad \forall t \in \interval[open]{t_n}{t_n + \delta}, \,\forall n \in \mathbb{N}.
  \end{equation}
  Combinando \eqref{eq:barbalat_contr}, \eqref{eq:barbalat_uc2} e \eqref{eq:barbalat_triang2} si ottiene
  \begin{equation*}
    \mleft| \dot{f}(t) \mright| \geq \mleft| \dot{f}(t_n) \mright| - \mleft| \dot{f}(t) - \dot{f}(t_n) \mright| > \varepsilon - \frac{\varepsilon}{2} = \frac{\varepsilon}{2}
  \end{equation*}
  da cui segue che
  \begin{equation}\label{eq:barbalat_contr3}
    \mleft| \dot{f}(t) \mright| > \frac{\varepsilon}{2} \quad \forall t \in \interval[open]{t_n}{t_n + \delta}, \,\forall n \in \mathbb{N}.
  \end{equation}
  Per ipotesi, per ogni $n \in \mathbb{N}$, $f$ è continua su $\interval{t_n}{t_n + \delta}$ e derivabile su $\interval[open]{t_n}{t_n + \delta}$, dunque, per il \thmcita{Teorema di Lagrange} e la relazione \eqref{eq:barbalat_contr3}, esiste $\xi_n \in \interval[open]{t_n}{t_n + \delta}$ tale che
  \begin{equation*}
    \mleft|f(t_n + \delta) - f(t_n)\mright| = \mleft|\dot{f}(\xi_n) \cdot \delta \mright| = \mleft|\dot{f}(\xi_n)\mright| \delta > \frac{\varepsilon\delta}{2}
  \end{equation*}
  da cui segue che
  \begin{equation}\label{eq:barbalat_contr4}
    \mleft|f(t_n + \delta) - f(t_n)\mright| > \frac{\varepsilon\delta}{2} > 0 \quad \forall n \in \mathbb{N}.
  \end{equation}
  Per ipotesi
  \[
  \lim_{n \to +\infty} f(t_n + \delta) - f(t_n) = \mathscr{l} - \mathscr{l} = 0,
  \]
  applicando la definizione di limite:
  \begin{equation*}
    \forall \sigma > 0, \, \exists \nu \in \mathbb{N} \, \colon \, \mleft|f(t_n + \delta) - f(t_n)\mright| < \sigma \quad \forall n > \nu. 
  \end{equation*}
  Scegliendo $\sigma = \varepsilon\delta/2$ si ottiene la relazione
  \begin{equation*}
    \mleft|f(t_n + \delta) - f(t_n)\mright| < \frac{\varepsilon\delta}{2} \quad \forall n > \nu
  \end{equation*}
  che è in contraddizione con \eqref{eq:barbalat_contr4}. L'assurdo nasce dall'aver negato la tesi, dunque, necessariamente, 
  \[
  \lim_{t \to + \infty} \dot{f}(t) = 0.
  \]
\end{proof}

\medskip
Se la funzione $\dot{f}$ non è uniformemente continua, non c'è garanzia che la tesi del \thmcita{Lemma di Barbalat \ref{lemma:barbalat}} sia vera. Viene mostrato, di seguito, un controesempio a sostegno di questa affermazione.

\section{Analisi del controesempio}
\begin{example}
Si consideri la funzione
\begin{equation}
    f(t) = \frac{\sin t^2}{t}
\end{equation}
definita su $\interval[open right]{0}{+\infty}$. Per il \thmcita{Teorema dei Carabinieri} si osserva subito che
\begin{equation*}
    \lim_{t \to +\infty} f(t) = 0 \,.
\end{equation*}
Calcolando la derivata prima si ottiene:
\begin{equation*}
    \dot{f}(t) = \diff*{\mleft(\frac{\sin t^2}{t}\mright)}{t} = \frac{2t \cos t^2 \cdot t - \sin t^2}{t^2} = 2\cos t^2 - \frac{\sin t^2}{t^2}
\end{equation*}

\smallskip
Siano definite le seguenti successioni:
\begin{equation*}
    x_n = \SQRT{2n\pi}, \quad y_n = \SQRT{2n\pi + \frac{\pi}{2}} \,.
\end{equation*}
La successione definita come $x_n - y_n$ è della forma
\begin{equation*}
    x_n - y_n = \SQRT{2n\pi} - \SQRT{2n\pi + \frac{\pi}{2}} = \frac{-\frac{\pi}{2}}{\SQRT{2n\pi} + \SQRT{2n\pi + \frac{\pi}{2}}} \,,
\end{equation*}
da cui segue che
\begin{equation}\label{eq:barbalat_contr5}
    \lim_{n \to +\infty} \mleft| x_n - y_n \mright| = 0.
\end{equation}
Si osserva che, per ogni $n \in \mathbb{N}$, $\dot{f}$ valutata in $x_n$ vale
\begin{equation*}
    \dot{f}(x_n) = 2\cos(2n\pi) - \frac{\sin(2n\pi)}{2n\pi} = 2(1) - 0 = 2 \, ,
\end{equation*}
mentre, per ogni $n \in \mathbb{N}$, $\dot{f}$ valutata in $y_n$ vale
\begin{equation*}
    \dot{f}(y_n) = 2\cos\mleft(2n\pi + \frac{\pi}{2}\mright) - \frac{\sin\mleft(2n\pi + \frac{\pi}{2}\mright)}{2n\pi + \frac{\pi}{2}} = 2(0) - \frac{1}{2n\pi + \frac{\pi}{2}} = - \frac{1}{2n\pi + \frac{\pi}{2}} \,.
\end{equation*}
Dunque, 
\begin{equation*}
  \mleft| \dot{f}(x_n) - \dot{f}(y_n) \mright| > 2 \quad \forall n \in \mathbb{N}
\end{equation*}
e, in particolare,
\begin{equation*}
    \lim_{n \to +\infty} \mleft| \dot{f}(x_n) - \dot{f}(y_n) \mright| = \lim_{n \to +\infty} \mleft| 2 + \frac{1}{2n\pi + \frac{\pi}{2}} \mright| = 2 \,.
\end{equation*}

Si scelga $\varepsilon = 2 > 0$ e si fissi un arbitrario $\delta > 0$. 
Dalla \eqref{eq:barbalat_contr5} segue che
\begin{equation*}
    \exists \nu \in \mathbb{N} \, \colon \, \mleft| x_n - y_n \mright| < \delta \quad \forall n > \nu.
\end{equation*}
Fissando un qualsiasi indice $\bar{n} > \nu$, si ottiene la relazione
\begin{equation*}
    \mleft| x_{\bar{n}} - y_{\bar{n}} \mright| < \delta.
\end{equation*}
Tuttavia, valutando la distanza tra le rispettive immagini, si osserva che
\begin{equation*}
    \mleft| \dot{f}(x_{\bar{n}}) - \dot{f}(y_{\bar{n}}) \mright| = 2 + \frac{1}{2\bar{n}\pi + \frac{\pi}{2}} > 2 = \varepsilon \,.
\end{equation*}
Si è così esibita una coppia di punti che, pur distando meno dell'arbitraria tolleranza $\delta$, possiedono immagini che distano strettamente più di $\varepsilon$. Essendo ciò costruttivamente valido per ogni possibile scelta di $\delta > 0$, la definizione di uniforme continuità per $\dot{f}$ risulta formalmente violata.

\smallskip
Infine, le medesime successioni $x_n$ e $y_n$ permettono di dimostrare in modo immediato che $\dot{f}$ non ammette limite per $t \to +\infty$. Si osserva banalmente che
\begin{equation*}
    \lim_{n \to +\infty} x_n = +\infty \quad \text{e} \quad \lim_{n \to +\infty} y_n = +\infty \,,
\end{equation*}
ma, passando al limite le rispettive immagini, si ottiene:
\begin{align*}
    \lim_{n \to +\infty} \dot{f}(x_n) &= \lim_{n \to +\infty} 2 = 2 \\
    \lim_{n \to +\infty} \dot{f}(y_n) &= \lim_{n \to +\infty} \mleft( - \frac{1}{2n\pi + \frac{\pi}{2}} \mright) = 0 \,
\end{align*}
da cui si deduce necessariamente che
\begin{equation*}
    \nexists \lim_{t \to +\infty} \dot{f}(t) \,.
\end{equation*}
\end{example}
\end{document}